% Generated draft: requires mathematical and source review.
\documentclass[11pt]{article}
\usepackage[margin=25mm]{geometry}
\usepackage{fontspec}
\setmainfont{Latin Modern Roman}
\usepackage{amsmath,amssymb,mathtools,unicode-math}
\setmathfont{Latin Modern Math}
\usepackage{enumitem,needspace,xurl,hyperref}
\setlength{\parindent}{0pt}
\setlength{\parskip}{5pt}
\setlength{\emergencystretch}{3em}
\newcommand{\N}{\mathbb{N}}
\newcommand{\Z}{\mathbb{Z}}
\newcommand{\Q}{\mathbb{Q}}
\newcommand{\R}{\mathbb{R}}
\newcommand{\C}{\mathbb{C}}
\newcommand{\F}{\mathbb{F}}
\newcommand{\A}{\mathbb{A}}
\newcommand{\PP}{\mathbb{P}}
\newcommand{\E}{\mathbb{E}}
\newcommand{\eps}{\varepsilon}
\newcommand{\dd}{\,\mathrm{d}}
\begin{document}
\section*{UnsolvedMath ID: 30001671}
\textbf{Source problem code:} catalogue entry
\subsection{Definitions and mathematical statement}
For each positive integer \(n\), let \([n]=\{1,\ldots,n\}\), and identify \(S_n\) with the set of words of length \(n\) in which every symbol of \([n]\) occurs exactly once. A word over \([n]\) is a finite sequence of symbols from \([n]\); concatenation \(xy\) places the sequence \(y\) after the sequence \(x\). Let \([n]^*\) denote all such words, including the empty word, and let \([n]^+\) denote the nonempty words. A \(k\)-colouring is a function \(c:S_n\to\{1,\ldots,k\}\), with no requirement that every colour occur.

The supplied statement leaves unresolved whether \(x,y,z\) must be nonempty. To record this qualification explicitly, choose either \(\mathcal W_n=[n]^*\) for every \(n\), or \(\mathcal W_n=[n]^+\) for every \(n\). For that choice, the proposed assertion is
\[
\begin{gathered}
\forall k\in\Z_{\geq1}\ \exists n\in\Z_{\geq1}\ 
\forall c:S_n\to\{1,\ldots,k\}\ 
\exists x,y,z\in\mathcal W_n:\\
xyz,xzy,yxz,yzx,zxy,zyx\in S_n,\\
c(xyz)=c(xzy)=c(yxz)=c(yzx)=c(zxy)=c(zyx).
\end{gathered}
\]
Membership of the concatenations in \(S_n\) requires that the letters occurring in \(x,y,z\), taken together, contain each element of \([n]\) exactly once. The source does not separately require the six displayed words to be distinct. Allowing empty component words permits coincidences among them. The nonempty-word version is therefore a possible intended formulation, but its nonemptiness hypothesis is not explicit in the supplied record and requires confirmation. No choice between these formulations is asserted here.
\subsection{Short English statement}
For every positive number of colours, is there a permutation size such that every colouring contains three blocks whose six possible concatenation orders are permutations of that size and all receive the same colour? The source does not specify whether the blocks must be nonempty; this qualification remains unresolved.
\subsection{Sources}
\begin{itemize}
\item Supplied local UnsolvedMath registry record, permanent ID \(30001671\), title \emph{Monochromatic Concatenations of Permutation Words}, source problem code \texttt{catalogue entry}, in the collection \emph{Oberwolfach Reports Open Problems}.
\item Supplied report citation: \emph{Combinatorics} (2011), Oberwolfach Reports, pp.\ 5--83, DOI \texttt{10.4171/owr/2011/01}. Source URL: \url{https://doi.org/10.4171/owr/2011/01}. The supplied background describes Oberwolfach Reports as documenting mathematical workshops and their open-problem discussions. No individual proposer is identified; the supplied proposed year is 2011.
\item Supplied literature reference, labelled ``Oberwolfach problem statement'': \url{https://doi.org/10.4171/owr/2011/01}.
\item Literature assessment dated exactly \texttt{2026-08-21}: the supplied curation labels the problem open and reports that no proof or counterexample for the exact monochromatic six-concatenation assertion was found. This is supplied dated curation derived from the OWR material, not an independently verified literature assessment. The supplied background requests independent verification of time-sensitive status and cited claims.
\item The supplied statement status is \texttt{corrected\_verified}; its accompanying verification note says only that TeX notation was enclosed in mathematical delimiters without changing mathematical content. This label does not establish mathematical review of the formulation above.
\end{itemize}
\end{document}
